% Generated by build/make_tables.py from gs_var_average_daily.csv. Do not edit.

\begin{table}[tbp]
\centering
\caption{Goldman Sachs average daily Value-at-Risk by risk category, and the disclosed diversification effect, FY2016--FY2025}
\label{tab:disclosed}
\begin{tabular}{@{}lrrrrrrrr@{}}
\toprule
 & \multicolumn{4}{c}{Risk category} & & & & \\
\cmidrule(lr){2-5}
Year & Rates & Equities & FX & Commod. & Summed & Diversif. & Total & \% of sum \\
\midrule
2016 & 45 & 25 & 21 & 17 & 108 & -45 & 63 & 41.7 \\
2017 & 40 & 24 & 12 & 13 & 89 & -35 & 54 & 39.3 \\
2018 & 46 & 31 & 14 & 11 & 102 & -42 & 60 & 41.2 \\
2019 & 46 & 27 & 11 & 12 & 96 & -40 & 56 & 41.7 \\
2020 & 71 & 55 & 23 & 20 & 169 & -75 & 94 & 44.4 \\
2021 & 60 & 43 & 13 & 25 & 141 & -55 & 86 & 39.0 \\
2022 & 97 & 33 & 32 & 47 & 209 & -95 & 114 & 45.5 \\
2023 & 96 & 29 & 24 & 19 & 168 & -69 & 99 & 41.1 \\
2024 & 81 & 37 & 26 & 19 & 163 & -71 & 92 & 43.6 \\
2025 & 70 & 47 & 24 & 17 & 158 & -68 & 90 & 43.0 \\
\bottomrule
\end{tabular}
\notetext{In millions of US dollars, except the final column. Average daily 95\% one-day Value-at-Risk as reported in the Market Risk Management section of each fiscal year's Form 10-K, taken from the filing that first reported the year rather than from a later filing's restatement. \emph{Summed} is the sum of the four category figures and is not itself disclosed. The diversification effect is disclosed as a negative number and reproduced with its sign; the final column expresses its magnitude as a percentage of the summed categories. Total VaR equals the summed categories plus the diversification effect in every year, exactly. The percentage averages 42.0 with a standard deviation of 2.08 across the ten years.}
\end{table}
